\chapter{Otros métodos de resolución}
\section{La transformada de Laplace}
Dada
\[x^{\left(n\right)} + a_{n-1}x^{\left(n-1\right)} + \cdots + a_{1}x' + a_{0}x = b\left(t\right), \]
sabemos que la podemos resolver si $\displaystyle b\left(t\right) $ es buena, es decir, por lo menos continua. Si no lo es, como es habitual en el contexto físico, no podemos obtener una solución particular para el completo. Por ejemplo, si $\displaystyle b\left(t\right) $ es continua a trozos, tanto el método de variación de constantes como el de parámetros indefinidos no funciona. 
\begin{eg}
Consideremos 
\end{eg}
\begin{definition}[Transformada de Laplace]
Dada $\displaystyle f : \R \to \R $, se define la \textbf{transformada de Laplace} de la forma
\[\mathcal{L}\left(f\left(t\right)\right) : = \int^{\infty}_{0} e^{-st}f\left(t\right) \; dt .\]
\end{definition}
\begin{definition}[Función continua a trozos]
	Decimos que $\displaystyle f : [a,b) \to \R $ es \textbf{continua a trozos} en $\displaystyle \left[a,b\right]  $ si es continua en todo el intervalo $\displaystyle \left[a,b\right]  $ salvo en un número finito de puntos en los que sí existe los límites laterales.
\end{definition}
\begin{definition}[Orden exponencial]
Decimos que $\displaystyle f : \left(0,\infty\right)\to \R $ es de \textbf{orden exponencial} si y solo si existen $\displaystyle K $, $\displaystyle M $ y $\displaystyle a $ tales que 
\[ \left|f\left(t\right)\right|\leq Ke^{at} ,\]
cuando $\displaystyle t \geq M $. 
\end{definition}
\begin{prop}[Condición suficiente de existencia]
Si $\displaystyle f : \left(0,\infty\right)\to \R $ es continua a trozos y de orden exponencial, entonces $\displaystyle \mathcal{L}\left(f\left(t\right)\right)<\infty $ para $\displaystyle s > a $. 
\end{prop}
\begin{proof}
Es sencillo ver que 
\[\mathcal{L}\left(f\left(t\right)\right)=\int^{\infty}_{0} e^{-st}f\left(t\right) \; dt = \int^{M}_{0} e^{-st}f\left(t\right) \; dt + \int^{\infty}_{0} e^{-st}f\left(t\right) \; dt .\]
La primera integral existe claramente y la segunda también puesto que
\[ \left|f\left(t\right)\right|\leq Ke^{at}.\]
\end{proof}
\begin{eg}
Consideremos la función 
\[f\left(t\right)=
\begin{cases}
3e^{t ^{2}}, \; t < 5 \\
10^{3}, \; t \geq 5
\end{cases}
.\]
Por la proposición anterior es fácil ver que la transformada de Laplace existe. También lo podemos ver de la forma:
\[\mathcal{L}\left(f\left(t\right)\right)=\int^{5}_{0} 3e^{t ^{2}}e^{-st} \; dt + \int^{\infty}_{5} 10^{3}e^{-st} \; dt < \infty .\]
\end{eg}
\begin{observation}
Nos preguntamos, si $\displaystyle \mathcal{L}\left(f\left(t\right)\right)=\mathcal{L}\left(g\left(t\right)\right) $, entonces $\displaystyle f\left(t\right) = g\left(t\right) $. Tenemos que por la linealidad de $\displaystyle \mathcal{L} $:
\[\mathcal{L}\left(f\left(t\right)\right)-\mathcal{L}\left(g\left(t\right)\right)=0 \iff \mathcal{L}\left(f\left(t\right)-g\left(t\right)\right)=0 \iff \int^{\infty}_{0} e^{st}\left(f\left(t\right)-g\left(t\right)\right) \; dt = 0 .\]
Como $\displaystyle e^{st} > 0 $, debe ser que $\displaystyle f\left(t\right)-g\left(t\right)=0 $ salvo quizás en un conjunto de medida nula.
\end{observation}
\begin{theorem}[Teorema de Lerch]
Sean $\displaystyle f $ y $\displaystyle g $ funciones de orden exponencial y continuas a trozos tales que $\displaystyle \mathcal{L}\left(f\left(t\right)\right)=\mathcal{L}\left(g\left(t\right)\right) $. Entonces, $\displaystyle f\left(t\right)=g\left(t\right) $, $\displaystyle \forall t > 0 $ salvo en sus respectivos puntos de discontinuidad.
\end{theorem}
Podemos entonces asegurar la existencia de $\displaystyle \mathcal{L}^{-1} $ en un cierto rango y nocvamos a ver la expresión explícita de $\displaystyle \mathcal{L}^{-1} $ porque es compleja, pero trabajaremos con ella a través de sus propiedades.
\begin{prop}
$\displaystyle \mathcal{L}^{-1} $ es lineal.
\end{prop}
\begin{proof}
Sean $\displaystyle \alpha, \beta \in \R $, entonces queremos ver que
\[
	\mathcal{L}^{-1}\left(\alpha F\left(s\right)+\beta G\left(s\right)\right)=  \alpha\mathcal{L}^{-1}\left(F\left(s\right)\right) + \beta\mathcal{L}^{-1}\left(G\left(s\right)\right).
\]
Si tomamos $\displaystyle F\left(s\right)=\mathcal{L}\left(f\left(t\right)\right) $ y $\displaystyle G\left(s\right)=\mathcal{L}\left(g\left(t\right)\right) $, tenemos que 
\[
\begin{split}
	\mathcal{L}^{-1}\left(\alpha F\left(s\right)+\beta G\left(s\right)\right)=& \mathcal{L}^{-1}\left(\alpha\mathcal{L}\left(f\left(t\right)\right)+\beta\mathcal{L}\left(g\left(t\right)\right)\right) = \mathcal{L}^{-1}\left(\mathcal{L}\left(\alpha f\left(t\right)+\beta g\left(t\right)\right)\right) \\
	= & \alpha f\left(t\right)+\beta g\left(t\right) = \alpha \mathcal{L}^{-1}\left(F\left(s\right)\right)+\beta \mathcal{L}^{-1}\left(G\left(s\right)\right) .
\end{split}
\]
\end{proof}
\begin{eg}
Sabiendo que para $\displaystyle f\left(t\right)=\sin kt $ tenemos que $\displaystyle F\left(s\right)=\frac{k}{s^{2}+k^{2}} $, buscamos $\displaystyle \mathcal{L}^{-1}\left(\frac{1}{s^{2}+7}\right) $. Es fácil ver que $\displaystyle f\left(t\right)=\frac{\sin7t}{\sqrt{7}} $.  
\end{eg}
\begin{eg}
Calculemos $\displaystyle f\left(t\right) $ tal que $\displaystyle F\left(s\right)=\frac{-2s + 6}{s^{2}+4} $. Tendremos que 
\[\mathcal{L}^{-1} \left(\frac{-2s + 6}{s^{2}+4}\right) = - 2\mathcal{L}^{-1}\left(\frac{s}{s^{2}+4}\right)+3\mathcal{L}^{-1}\left(\frac{2}{s^{2}+4}\right)=-2\cos2t+3\sin2t.\]
\end{eg}
\begin{prop}[Cambio de escala]
Si $\displaystyle a > 0 $, tenemos que 
\[\mathcal{L}\left(f\left(at\right)\right)=\int^{\infty}_{0} e^{-st}f\left(at\right) \; dt = \frac{1}{a}\int^{\infty}_{0} e^{-\frac{s}{a}u}f\left(u\right) \; du = \frac{1}{a}F\left(\frac{s}{a}\right) .\]
\end{prop}
\begin{observation}[Traslaciones]
	Podemos hacer un par de observaciones.
	\begin{itemize}
	\item Para cualquier $\displaystyle a \in \R $, tenemos que $\displaystyle \mathcal{L}\left(e^{at}f\left(t\right)\right)=F\left(s-a\right) $.
	\item Para $\displaystyle a > 0 $, $\displaystyle \mathcal{L}\left(H_{a}\left(t\right)f\left(t-a\right)\right)=e^{-as}F\left(s\right) $. 
	\end{itemize}
	En efecto,
	\[\mathcal{L}\left(e^{at}f\left(t\right)\right)=\int^{\infty}_{0} e^{at}e^{-st}f\left(t\right) \; dt = \int^{\infty}_{0} e^{-\left(s-a\right)t}f\left(t\right) \; dt=F\left(s-a\right) .\]
Por otro lado,
\[\mathcal{L}\left(H_{a}\left(t\right)f\left(t-a\right)\right)=\int^{\infty}_{0} e^{-st}f\left(t-a\right) \; dt=\int^{\infty}_{a} e^{-sa}e^{-s\left(t-a\right)}f\left(t-a\right) \; dt=e^{-sa}\int^{\infty}_{0} e^{-su}f\left(u\right) \; du=e^{-sa}F\left(s\right) .\]
\end{observation}
\begin{prop}[Derivación]
	Para $\displaystyle f\left(t\right) $ con $\displaystyle \mathcal{L}\left(f\left(t\right)\right)=F\left(s\right) $ tenemos que 

\end{prop}

